%% ==============================================================================
%% content/05-conclusion.tex
%% Chapitre 5 : Conclusion Générale et Perspectives
%% Encadre la thèse symétriquement à l'Introduction Générale
%% ==============================================================================
\chapter{Conclusion Générale et Perspectives}
\label{chap:conclusion-generale}

\section{Le Statut de la Conclusion Générale en Thèse}
\label{sec:conclusion-statut}

Dans la structure académique d'une thèse de doctorat ou d'un mémoire d'habilitation,
l'introduction générale et la conclusion générale occupent une place singulière.
Alors que les chapitres intermédiaires développent les contributions scientifiques techniques pas à pas,
la conclusion générale prend de la hauteur pour dresser le bilan critique des travaux,
évaluer la portée des résultats obtenus et tracer les lignes de fuite pour les recherches futures.

Noyer la conclusion à la fin du dernier chapitre de résultats constitue un écueil fréquent.
En isolant la conclusion générale dans son propre module (\code{content/05-conclusion.tex} ou \code{content/conclusion.tex}),
vous offrez au jury et aux lecteurs une vision d'ensemble claire et détendue des résultats.

\begin{callout}[Numérotation : \code{\textbackslash chapter} ou \code{\textbackslash addchap} ?]
Selon les traditions doctorales de votre établissement, deux conventions sont possibles :
\begin{enumerate}
  \item \textbf{Chapitre numéroté classique (\code{\textbackslash chapter\{Conclusion Générale\}}) :}
  s'intègre dans le décompte séquentiel des chapitres (ex. Chapitre~5).
  \item \textbf{Chapitre d'ouverture/fermeture non numéroté (\code{\textbackslash addchap\{Conclusion Générale\}}) :}
  KOMA-Script propose la commande dédiée \code{\textbackslash addchap}.
  Elle compose le titre avec la même majesté typographique qu'un chapitre numéroté,
  l'insère dans la table des matières (\code{TOC})
  et met à jour les en-têtes de page courants, tout en omettant le préfixe numérique.
  Cette approche permet d'encadrer symétriquement le corps d'ouvrage
  entre \code{\textbackslash addchap\{Introduction Générale\}} et \code{\textbackslash addchap\{Conclusion Générale\}}.
\end{enumerate}
\end{callout}

\section{Synthèse des Apports de l'Architecture Modulaire}
\label{sec:conclusion-synthese}

Le canevas développé tout au long de ce tutoriel résout les limitations historiques
des projets \LaTeX{} volumineux en formalisant six piliers méthodologiques :

\begin{enumerate}
  \item \textbf{Le découplage sémantique total :}
  le fichier maître \code{main.tex} ordonnance la structure sans être pollué par le style ;
  la classe \code{mythesis.cls} encapsule l'identité visuelle KOMA-Script ;
  et les chapitres rédactionnels (\code{content/*.tex}) se concentrent exclusivement sur le fond.
  \item \textbf{La flexibilité du support par fichiers d'options (\code{.clo}) :}
  le passage d'un format de lecture sur écran (marges symétriques, palette interactive)
  à un format d'impression papier reliée (recto-verso, décalage de pliure \code{BCOR=10mm}, liens monochromes)
  s'effectue par la simple modification d'un mot-clé dans \code{main.tex}.
  \item \textbf{La sanctuarisation des données brutes et du code source :}
  l'introduction des dossiers dédiés \code{data/} (pour les jeux de données tabulaires \code{.csv} lus par \code{pgfplots})
  et \code{listings/} (pour les codes sources externes importés sans copier-coller)
  garantit une traçabilité expérimentale irréprochable.
  \item \textbf{La reproductibilité typographique multi-plateforme :}
  le stockage local des polices OpenType/TrueType dans le répertoire \code{fonts/}
  immunise le projet contre les variations de polices entre Windows, Linux, macOS, Overleaf et les pipelines d'intégration continue CI/CD.
  \item \textbf{La bibliographie modulaire Bib\LaTeX{} :}
  la séparation entre données brutes (\code{references.bib}),
  mise en page de la bibliographie finale (\code{custom-style.bbx})
  et formatage des citations cliquables (\code{custom-style.cbx})
  offre une liberté totale de personnalisation sans la complexité des anciens fichiers BST.
  \item \textbf{L'outillage qualité automatisé :}
  la présence conjointe de \code{latexmkrc} (confinement de la compilation dans \code{build/}),
  de \code{.editorconfig} (uniformisation de l'indentation à deux espaces)
  et de \code{.chktexrc} (contrôle automatique des espaces insécables avant \code{\textbackslash cite} et \code{\textbackslash ref})
  élève le document aux standards industriels du génie logiciel.
\end{enumerate}

\section{Perspectives et Évolutions du Modèle}
\label{sec:conclusion-perspectives}

Cette architecture constitue un socle extensible qui ouvre plusieurs voies d'enrichissement :

\begin{itemize}
  \item \textbf{Intégration Continue (GitHub Actions / GitLab CI) :}
  l'exécution automatique de \code{latexmk} dans un conteneur TeX Live à chaque \code{git push}
  permet de générer et de publier les PDF d'évaluation (écran et papier) sous forme d'artefacts téléchargeables.
  \item \textbf{Variantes de profils d'options :}
  créer de nouveaux fichiers \code{.clo}, tels qu'un profil pour liseuse à encre électronique (\code{mythesis-ereader.clo})
  ou un profil accessible à grand corps de texte (\code{mythesis-dyslexic.clo}).
  \item \textbf{Extraction automatique de données :}
  connecter des scripts Python générant automatiquement les courbes de \code{data/}
  lors de l'exécution de bancs d'essai scientifiques.
\end{itemize}

En adoptant ces préceptes, vous transformez votre projet \LaTeX{} en une plateforme d'édition pérenne,
élégante et véritablement partagée.

